\documentclass[11pt]{amsart}
\usepackage[margin=1in]{geometry}
\usepackage{amssymb}
\newtheorem{theorem}{Theorem}
\newtheorem{corollary}[theorem]{Corollary}
\newtheorem{lemma}[theorem]{Lemma}
\theoremstyle{remark}
\newtheorem{remark}[theorem]{Remark}
\newcommand{\R}{\mathbb{R}}
\DeclareMathOperator{\re}{Re}
\begin{document}
\title[A Gaussian form of Weil's positivity criterion]{A Gaussian form of Weil's positivity criterion\\ for the Riemann hypothesis}
\author{James C. Kerr}
\date{November 19, 2026}
\subjclass[2020]{Primary 11M26; Secondary 11M06, 68V20}
\keywords{Riemann hypothesis, Weil's criterion, explicit formula, Gaussian test functions, formal verification}
\begin{abstract}
We prove that the Riemann hypothesis is equivalent to the nonnegativity of an explicit real function
$A_b(\alpha)$ of two real variables, built from the primes and the gamma function, for all $\alpha\in\R$ and $0<b\le 1$.
By the explicit formula $A_b(\alpha)$ is the Gaussian-smoothed density of the zeta zeros, so the result
restricts Weil's positivity criterion to a two-parameter family of Gaussian test functions.
Two corollaries follow: a two-point log-convexity criterion, and a criterion in which only a sum over
primes appears. All statements have been formally verified in Lean~4 with Mathlib.
\end{abstract}
\maketitle

\section{Introduction}
Write each nontrivial zero of $\zeta$ as $\rho=\tfrac12+\eta+i\gamma$ with $\eta,\gamma\in\R$,
$|\eta|<\tfrac12$, and put $\tau_\rho=\gamma-i\eta$, so that $\rho=\tfrac12+i\tau_\rho$. The Riemann
hypothesis (RH) says that every $\tau_\rho$ is real. For $b>0$ and $\alpha\in\R$ let
\[
g_{b,\alpha}(u)=e^{-bu^2-i\alpha u},\qquad
G_{b,\alpha}(\tau)=\int_{\R}g_{b,\alpha}(u)e^{i\tau u}\,du=\sqrt{\pi/b}\;e^{-(\tau-\alpha)^2/(4b)}\quad(\tau\in\mathbb{C}),
\]
and define, with $\psi=\Gamma'/\Gamma$ and $\Lambda$ the von Mangoldt function,
\begin{align*}
P_b(\alpha)&=2\sqrt{\pi/b}\;e^{(1-4\alpha^2)/(16b)}\cos\!\big(\alpha/(4b)\big),\\
W_b(\alpha)&=\frac{1}{\sqrt\pi}\int_{\R}e^{-u^2}\Big(\re\psi\big(\tfrac14+\tfrac i2(\alpha+2\sqrt b\,u)\big)-\log\pi\Big)\,du,\\
S_b(\alpha)&=2\sum_{n\ge2}\frac{\Lambda(n)}{\sqrt n}\,e^{-b\log^2 n}\cos(\alpha\log n),\qquad
A_b(\alpha)=P_b(\alpha)+W_b(\alpha)-S_b(\alpha).
\end{align*}
The function $A_b(\alpha)$ involves no zeros. Weil's explicit formula \cite{Weil}, in the form of Barner
\cite{Barner} (see also \cite[Ch.~XVII]{Lang}), applies to $g_{b,\alpha}$, which is smooth with
$g_{b,\alpha},g_{b,\alpha}'=O(e^{-|u|})$, and gives
\begin{equation}\label{eq:EF}
\sum_\rho G_{b,\alpha}(\tau_\rho)=A_b(\alpha)\qquad(b>0,\ \alpha\in\R),
\end{equation}
the sum running over the nontrivial zeros with multiplicity and converging absolutely. Thus $A_b(\alpha)$
is the counting measure of the $\tau_\rho$ smoothed by a Gaussian of variance $2b$ centred at $\alpha$.

\begin{theorem}\label{thm:main}
The Riemann hypothesis holds if and only if $A_b(\alpha)\ge0$ for all $\alpha\in\R$ and all $0<b\le1$.
\end{theorem}

\begin{corollary}\label{cor:two}
The Riemann hypothesis holds if and only if for every $0<b\le1$ there is $\alpha_0$ with $A_b(\alpha_0)>0$ and
\begin{equation}\label{eq:two}
A_b(\alpha+h)^2\le e^{h^2/(2b)}A_b(\alpha)\,A_b(\alpha+2h)\qquad(\alpha,h\in\R).
\end{equation}
\end{corollary}

\begin{theorem}\label{thm:primes}
Let $T_b(\alpha)=2\sum_{p}p^{-1/2}\log p\;e^{-b\log^2p}\cos(\alpha\log p)$, the sum over primes. The Riemann
hypothesis holds if and only if for every $\alpha\in\R$ there is $C(\alpha)$ such that
\begin{equation}\label{eq:primes}
\big|P_b(\alpha)-T_b(\alpha)\big|\le C(\alpha)\,b^{-1/2}\qquad(0<b\le1).
\end{equation}
\end{theorem}

\subsection*{Relation to known results}
Weil \cite{Weil} showed that RH is equivalent to the nonnegativity of the zero-side functional
$h\mapsto\sum_\rho\int_\R h(u)e^{i\tau_\rho u}\,du$ on autocorrelations $h=f*\tilde f$,
$\tilde f(u)=\overline{f(-u)}$, for $f$ in a suitable class of test functions. Since $g_{b,\alpha}=f*\tilde f$ with $f(u)=(4b/\pi)^{1/4}e^{-2bu^2-i\alpha u}$, the
condition in Theorem~\ref{thm:main} is Weil's condition for Gaussian $f$ only; the theorem says that this
two-parameter family already suffices. Gaussians belong to Weil's class, so necessity is a case of his lemma.
For sufficiency Weil takes, at an off-line zero $\tfrac12+\eta_0+i\gamma_0$, the test function with transform
$|P(\tau-\gamma_0)|^2e^{-2A(\tau-\gamma_0)^2}$ on the line, where $P(z)=zQ(z)Q(-z)$: the factor $z$ makes the term of
the off-line zero negative and $Q$ vanishes at the other nearby zeros \cite[p.~263]{Weil}; a Lean formalization
of this argument is given in \cite{Lin}. Theorem~\ref{thm:main} shows that no polynomial factor is needed: the
sign comes from the phase $e^{is\theta_\alpha(\rho)}$ after moving the centre $\alpha$ off $\gamma_0$. Sufficiency
can also be deduced from Weil's argument, since convolving his function with $G_{b,x}$ and using \eqref{eq:EF}
expresses its zero sum as a limit of integrals of $A_b(x)$ against a nonnegative weight; the direct proof below
gives in addition the rate in Remark~\ref{rem:margin}. Other restrictions of the test class
have been studied, notably to functions of compact support \cite{Yoshida,Bombieri} and to the sequence
underlying Li's criterion \cite{Li,BL}; Weil himself later asked for classes of functions on which the
positivity could be examined \cite[Commentaire]{Weil}. Related criteria with Gaussian test functions appear
in \cite{Dow}, for all even Schwartz functions against a Gaussian weight, and in \cite{Takahashi}, for translated
Gaussian windows on the critical line. For a reformulation of RH in a similar spirit see \cite{SD}.

\section{Proof of Theorem \ref{thm:main}}
Throughout $s=1/(4b)$, so $0<b\le1$ means $s\ge\tfrac14$. For a zero $\rho=\tfrac12+\eta+i\gamma$ and $\alpha\in\R$ put
\[
v_\alpha(\rho)=\eta^2-(\gamma-\alpha)^2,\qquad\theta_\alpha(\rho)=2\eta(\gamma-\alpha).
\]
Since $-(\tau_\rho-\alpha)^2=v_\alpha(\rho)+i\theta_\alpha(\rho)$, we have
$G_{b,\alpha}(\tau_\rho)=\sqrt{\pi/b}\,e^{sv_\alpha(\rho)}e^{is\theta_\alpha(\rho)}$. The zeros are symmetric under
$\rho\mapsto1-\bar\rho$, that is $(\eta,\gamma)\mapsto(-\eta,\gamma)$, with equal multiplicities, so by \eqref{eq:EF}
\begin{equation}\label{eq:real}
A_b(\alpha)=\sqrt{\pi/b}\,\sum_\rho e^{sv_\alpha(\rho)}\cos\big(s\,\theta_\alpha(\rho)\big).
\end{equation}

\begin{lemma}\label{lem:tail}
Fix $\alpha\in\R$ and let $Z_\alpha$ be the finite set of zeros with $|\gamma-\alpha|\le1$, of cardinality $|Z_\alpha|$ with multiplicity. There is $C=C(\alpha)$ with
$\sum_{\rho\notin Z_\alpha}e^{sv_\alpha(\rho)}\le C\,e^{-3s/4}$ for all $s\ge\tfrac14$.
\end{lemma}
\begin{proof}
Let $x=(\gamma-\alpha)^2>1$. As $\eta^2<\tfrac14$, $v_\alpha(\rho)<\tfrac14-x\le-\tfrac34x$, and for $s\ge\tfrac14$,
$\tfrac34sx\ge\tfrac34s+\tfrac3{16}(x-1)$. Hence $e^{sv_\alpha(\rho)}\le e^{-3s/4}e^{-\frac3{16}((\gamma-\alpha)^2-1)}$, and
$C=\sum_\rho e^{-\frac3{16}((\gamma-\alpha)^2-1)}$ is finite because the number of zeros with $T<\gamma\le T+1$ is
$O(\log(2+|T|))$ \cite[Thm.~9.2]{Titchmarsh}.
\end{proof}

\begin{lemma}\label{lem:dominant}
Suppose $\rho_0=\tfrac12+\eta_0+i\gamma_0$ is a zero with $\eta_0\ne0$. Then there are $\alpha\in\R$, an integer
$N\ge2$ and reals $M>0$, $\omega>0$, $\varepsilon>0$, $K\ge0$ such that for all $s\ge\tfrac14$
\begin{equation}\label{eq:dom}
\Big|\sum_\rho e^{sv_\alpha(\rho)}\cos\big(s\,\theta_\alpha(\rho)\big)-N e^{sM}\cos(s\omega)\Big|\le K\,e^{s(M-\varepsilon)} .
\end{equation}
\end{lemma}
\begin{proof}
Let $I=(\gamma_0-|\eta_0|,\gamma_0+|\eta_0|)$. For $\alpha\in I$ we have $v_\alpha(\rho_0)>0$. If $v_\alpha(\rho)>0$ for
some $\alpha\in I$ then $|\gamma-\alpha|<|\eta|<\tfrac12$, so $\eta\ne0$ and $|\gamma-\gamma_0|<1$; let $F$ be the finite
set of such zeros. For $\rho,\rho'\in F$ with $\gamma\ne\gamma'$ the difference
$v_\alpha(\rho)-v_\alpha(\rho')=\eta^2-\eta'^2-\gamma^2+\gamma'^2+2\alpha(\gamma-\gamma')$ is a nonconstant affine
function of $\alpha$ and vanishes for at most one $\alpha$. We may therefore fix $\alpha\in I$ such that
(a) $\alpha\ne\gamma$ for every $\rho\in F$, and (b) $v_\alpha(\rho)\ne v_\alpha(\rho')$ whenever $\rho,\rho'\in F$
have different ordinates. Let $M=\max_\rho v_\alpha(\rho)$; then $M\ge v_\alpha(\rho_0)>0$, and the set $D$ of zeros with
$v_\alpha(\rho)=M$ lies in $F$. By (b) all zeros in $D$ share one ordinate $\gamma_*$, hence one value of $\eta^2$, say
$\eta_*^2>0$; by the symmetry $\eta\mapsto-\eta$, $D$ consists of $\tfrac12\pm\eta_*+i\gamma_*$, and its cardinality
$N$ with multiplicity is at least $2$. By (a), $\omega:=2|\eta_*||\gamma_*-\alpha|>0$, and
$\cos(s\theta_\alpha(\rho))=\cos(s\omega)$ for $\rho\in D$. The zeros in $Z_\alpha\setminus D$ are finitely many and
satisfy $v_\alpha(\rho)\le M-\varepsilon$ for some $0<\varepsilon\le M+\tfrac34$. With $C$ from Lemma~\ref{lem:tail},
the left side of \eqref{eq:dom} is at most $|Z_\alpha|e^{s(M-\varepsilon)}+Ce^{-3s/4}\le(|Z_\alpha|+C)e^{s(M-\varepsilon)}$.
\end{proof}

\begin{proof}[Proof of Theorem \ref{thm:main}]
If RH holds, every $\eta=0$ and each term of \eqref{eq:real} equals $e^{-s(\gamma-\alpha)^2}>0$, so $A_b(\alpha)\ge0$.
Conversely, suppose $A_b(\alpha)\ge0$ for all $\alpha$ and $0<b\le1$, and that some zero has $\eta_0\ne0$. Take
$\alpha,N,M,\omega,\varepsilon,K$ from Lemma~\ref{lem:dominant} and $s_k=(2k+1)\pi/\omega$, $b_k=1/(4s_k)$. For large
$k$ we have $s_k\ge\tfrac14$, and by \eqref{eq:real} and \eqref{eq:dom}
\[
\sqrt{b_k/\pi}\;A_{b_k}(\alpha)\le-Ne^{s_kM}+Ke^{s_k(M-\varepsilon)}=-e^{s_kM}\big(N-Ke^{-\varepsilon s_k}\big)<0
\]
once $Ke^{-\varepsilon s_k}<N$, a contradiction. Hence every $\eta=0$.
\end{proof}

\begin{remark}\label{rem:margin}
The proof shows that an off-line zero forces $A_b(\alpha)\le-c\,b^{-1/2}e^{M/(4b)}$ along a sequence $b\to0$,
for some $\alpha$ and $c,M>0$. On the other hand, if $\gamma<\gamma'$ are consecutive ordinates of zeros on the
critical line with $g=\gamma'-\gamma$, the two nearest terms at $\alpha=(\gamma+\gamma')/2$ contribute
$2\sqrt{\pi/b}\,e^{-g^2/(16b)}$. Under RH the margin in Theorem~\ref{thm:main} is therefore exponentially small as
$b\to0$.
\end{remark}

\section{Proofs of Corollary \ref{cor:two} and Theorem \ref{thm:primes}}
\begin{proof}[Proof of Corollary \ref{cor:two}]
Assume RH. Then every $\tau_\rho=\gamma$ is real, and by \eqref{eq:real}
\[
F(\alpha):=\sqrt{b/\pi}\;e^{s\alpha^2}A_b(\alpha)=\sum_\gamma e^{-s\gamma^2}e^{2s\alpha\gamma}
\]
is a sum of positive exponentials, so $F(\alpha+h)^2\le F(\alpha)F(\alpha+2h)$ by the Cauchy--Schwarz inequality. Since
$\alpha^2+(\alpha+2h)^2-2(\alpha+h)^2=2h^2$ and $2sh^2=h^2/(2b)$, this is \eqref{eq:two}; and $A_b>0$ because
$\zeta$ has nontrivial zeros. Conversely, taking $\alpha+2h=\alpha_0$ in \eqref{eq:two} gives
$A_b(\alpha)A_b(\alpha_0)\ge0$, hence $A_b(\alpha)\ge0$ for every $\alpha$, and Theorem~\ref{thm:main} applies.
\end{proof}

\begin{proof}[Proof of Theorem \ref{thm:primes}]
Let $0<b\le1$. By Stirling's formula $|\re\psi(\tfrac14+\tfrac i2r)|\le c_0(1+\log(2+|r|))$, so
$|W_b(\alpha)|\le c_1(1+\log(2+|\alpha|))$. The prime powers $p^k$ with $k\ge3$ contribute to $S_b(\alpha)$ at most
$2\sum_p\log p\sum_{k\ge3}p^{-k/2}<\infty$ in absolute value. For $k=2$, with $w(x)=e^{-4b\log^2x}$ and Chebyshev's
bound $\vartheta(x)\le x\log4$, partial summation gives
\[
\sum_p\frac{\log p}{p}\,w(p)=-\int_1^\infty\vartheta(x)\,d\Big(\frac{w(x)}{x}\Big)
\le\log4\Big(w(1)+\int_1^\infty\frac{w(x)}{x}\,dx\Big)=\log4\Big(1+\frac{\sqrt\pi}{4\sqrt b}\Big).
\]
Hence $|A_b(\alpha)-(P_b(\alpha)-T_b(\alpha))|\le c_2(\alpha)\,b^{-1/2}$, and \eqref{eq:primes} holds for a given
$\alpha$ if and only if $r_s(\alpha):=\sum_\rho e^{sv_\alpha(\rho)}\cos(s\theta_\alpha(\rho))$ is bounded for
$s\ge\tfrac14$. Under RH, $0<r_s(\alpha)\le|Z_\alpha|+C(\alpha)$ by Lemma~\ref{lem:tail}. If RH fails, take $\alpha$
as in Lemma~\ref{lem:dominant} and $s_k=2k\pi/\omega$: then
$r_{s_k}(\alpha)\ge e^{s_kM}(N-Ke^{-\varepsilon s_k})\to\infty$, so \eqref{eq:primes} fails for this $\alpha$.
\end{proof}

\section{Formal verification}
\begin{sloppypar}
The statements above are checked in Lean~4 (v4.35.0) against Mathlib (v4.35.0) \cite{mathlib}. The
sources are available at \texttt{github.com/QuantQJ/weil-gaussian-lean}, and RH is Mathlib's
\texttt{RiemannHypothesis}. Each theorem below depends only on the standard axioms \texttt{propext},
\texttt{Classical.choice} and \texttt{Quot.sound}.
\end{sloppypar}

\smallskip
\begin{center}
\begin{tabular}{ll}
\eqref{eq:EF} for Gaussians & \texttt{FleetPrimePoleBoundRH.zero\_trace\_eq\_rhs}\\
Theorem \ref{thm:main}, zero side & \texttt{FleetHeatPositivityRH.RH\_iff\_heat\_positivity}\\
Theorem \ref{thm:main}, arithmetic side & \texttt{FleetHeatPositivityRH.RH\_iff\_rhs\_positivity}\\
Corollary \ref{cor:two} & \texttt{FleetTwoPointConvexity.RH\_iff\_twoPoint}\\
Theorem \ref{thm:primes} & \texttt{FleetPrimeOnlyRH.RH\_iff\_primeOnly\_bound}
\end{tabular}
\end{center}
\smallskip

\noindent In the formal statements $A_b(\alpha)$ is \texttt{(rhs (gauss b $\alpha$)).re}, where \texttt{rhs} is the sum
of the pole, archimedean and prime-power terms defined above. In \texttt{RH\_iff\_twoPoint} the existence of
$\alpha_0$ is a hypothesis of the equivalence; the remark that it holds under RH uses the classical existence of a
nontrivial zero and is not part of the formal statement. As a numerical check of \eqref{eq:EF},
$A_{1/2}(20)=1.486853331\ldots$ both from the prime powers below $4723$ and from the first $40$ zeros.

\subsection*{Acknowledgement}
The Lean development and this note were prepared with the assistance of AI systems; every formal statement is
checked by the Lean kernel.

\begin{thebibliography}{99}
\bibitem{Barner} K.~Barner, On A.~Weil's explicit formula, \emph{J. Reine Angew. Math.} \textbf{323} (1981), 139--152.
\bibitem{Bombieri} E.~Bombieri, Remarks on Weil's quadratic functional in the theory of prime numbers, I,
\emph{Rend. Mat. Acc. Lincei} (9) \textbf{11} (2000), 183--233.
\bibitem{BL} E.~Bombieri and J.~C.~Lagarias, Complements to Li's criterion for the Riemann hypothesis,
\emph{J. Number Theory} \textbf{77} (1999), 274--287.
\bibitem{Dow} J.~Dow, The Prime--Resonance Operator, I--II, preprints, Zenodo, 2025,
doi:10.5281/zenodo.17503095, doi:10.5281/zenodo.17605172.
\bibitem{Lang} S.~Lang, \emph{Algebraic Number Theory}, 2nd ed., Springer, 1994.
\bibitem{Li} X.-J.~Li, The positivity of a sequence of numbers and the Riemann hypothesis,
\emph{J. Number Theory} \textbf{65} (1997), 325--333.
\bibitem{Lin} C.-K.~Lin, Formalizing Weil's explicit formula and positivity criterion for the Riemann hypothesis
in Lean, preprint, Zenodo, 2026, doi:10.5281/zenodo.23494393.
\bibitem{mathlib} The mathlib Community, The Lean mathematical library, in \emph{Proc. CPP 2020}, ACM, 2020, 367--381.
\bibitem{SD} J.~Sondow and C.~Dumitrescu, A monotonicity property of Riemann's xi function and a reformulation
of the Riemann hypothesis, \emph{Period. Math. Hungar.} \textbf{60} (2010), 37--40.
\bibitem{Takahashi} T.~Takahashi, The WB-Bridge principle: refined spectral positivity and an RH-equivalent
criterion, preprint, Zenodo, 2025, doi:10.5281/zenodo.17574853.
\bibitem{Titchmarsh} E.~C.~Titchmarsh, \emph{The Theory of the Riemann Zeta-function}, 2nd ed., Oxford Univ. Press, 1986.
\bibitem{Weil} A.~Weil, Sur les ``formules explicites'' de la th\'eorie des nombres premiers,
\emph{Comm. S\'em. Math. Univ. Lund}, Tome Suppl. (1952), 252--265; reprinted with a Commentaire in
\emph{\OE uvres Scientifiques} II, Springer, 1979, 48--61, 527--528.
\bibitem{Yoshida} H.~Yoshida, On Hermitian forms attached to zeta functions, in \emph{Zeta Functions in Geometry},
Adv. Stud. Pure Math. \textbf{21}, 1992.
\end{thebibliography}
\end{document}
